\section{Introduction}

Since the first detection of a binary black-hole merger~\cite{GW150914} gravitational-wave detectors have produced a stream of strain data that is public through the Gravitational Wave Open Science Center
(GWOSC)~\cite{GWOSCO12}. Searches for \emph{short, unmodelled transients} (bursts) cannot rely on a bank of astrophysically motivated waveforms and must either use phenomenological
templates or make no assumption on the signal shape (excess-power and coherent-network methods~\cite{ExcessPower,cWB}). They all face the same statistical problem: the noise is only approximately
Gaussian and stationary, the signals are weak ($\snr\sim5$--$10$), the parameters of the signal are unknown nuisance parameters, and the required false-alarm probability is tiny.

This project uses that problem as a laboratory for the techniques of the course \emph{Statistics and Data Analysis}: hypothesis testing and the Neyman--Pearson lemma, maximum likelihood,
profile likelihood and Wilks' theorem, toy Monte Carlo, Bayesian inference with priors, MCMC and Bayes factors, supervised and unsupervised learning, calibration and uncertainty in machine learning. The organisation follows the
earlier projects on LHCb open data (asymmetry in $B^\pm\to K^\pm K^+K^-$) and on atomic-force-microscopy curves: one shared data set, one notebook per method, one common set of figures of merit, a report that says where each number comes from.

\paragraph{Research question.} \emph{At a fixed false-alarm probability on real LIGO noise, how do a frequentist matched-filter bank, a Bayesian analysis of the same likelihood and machine-learning detectors
compare in detection efficiency, in the quality of their uncertainty statements, and in robustness to the non-Gaussian noise and to signal morphologies they were not designed for? And, for a signal that is clearly above the noise, which structure can tell which morphology produced it?}

\paragraph{Scope.} Mostly a single detector (H1); three phenomenological burst morphologies, and binary-black-hole (IMRPhenomD) signals in Section~\ref{sec:imr}, injected at a chosen optimal SNR; false-alarm probabilities per one-second window; a second detector with coincidence and time slides only in Section~\ref{sec:coinc}. This is \emph{not} an astrophysical sensitivity estimate: SNRs are not converted into distances and the search is not a coherent network search.

\paragraph{Organisation.} Section~\ref{sec:data} describes the data, the whitening and the tests of the noise model. Section~\ref{sec:signals} defines the signals, the SNR and the injection chain. Sections~\ref{sec:freq}, \ref{sec:bayes} and \ref{sec:ml}
are the three analyses; Section~\ref{sec:seq} extends the machine-learning part with sequence models (Transformer, GRU), morphology classification (signal hypothesis testing) and a tuning study of the Transformer-against-CNN comparison; Section~\ref{sec:imr} repeats the detection with realistic binary-black-hole (IMRPhenomD) signals, Section~\ref{sec:coinc} adds a second detector with coincidence and time slides, and Section~\ref{sec:openset} tests how signals outside the hypotheses can be flagged; Section~\ref{sec:physical} expresses the sensitivities in physical units and Section~\ref{sec:pe} estimates the parameters of the binaries; Section~\ref{sec:unfold} unfolds their mass distribution; Section~\ref{sec:comparison} compares the approaches (with an ablation of the CNN advantage, the threshold convention and the generalisation to O2 and O3) and Section~\ref{sec:conclusions} gives conclusions and limits. Most sections end with subsections added in a second round of extensions: bootstrap errors and seeds, robustness, data-quality flags, a coherent two-detector statistic, unsupervised noise models and the limits of the one-second window. Appendix~\ref{app:code} maps every number to the notebook and result file that produced it.
